\section{IMPLEMENTATION \& SIMULATION}
\subsection{Tech Stack}
Game sử dụng những công cụ sau:
\begin{itemize}
	\item \textbf{Game Engine:} Godot Engine (Version 4.6.2) sử dụng ngôn ngữ lập trình GDScript.
	\item \textbf{Artificial Intelligence:} Mô hình ngôn ngữ lớn (LLM) Qwen2.5:3b, vận hành cục bộ thông qua nền tảng Ollama.
	\item \textbf{Database:} Hệ quản trị cơ sở dữ liệu quan hệ SQLite (tích hợp qua plugin Godot-SQLite).
	\item \textbf{Design \& Version Control:} Figma (Thiết kế UI/UX) và GitHub (Quản lý mã nguồn, đồng bộ làm việc nhóm).
\end{itemize}

\subsection{System Implementation}
\subsubsection*{a. Game Narrative Context (Bối cảnh Ứng dụng)}
Hệ thống kỹ thuật được xây dựng bám sát vào cốt truyện cốt lõi: Người chơi vào vai \textit{"The Seeker"} trong một thế giới đang bị hủy hoại bởi \textit{"Gibberish"} (thứ ngôn ngữ không quy tắc biến con người thành những cái xác rỗng \textbf{The Hollowed}). Tại Thánh Vực Tri Thức, thông qua sự hướng dẫn của thủ thư \textbf{Elaria (được mô phỏng bởi AI)}, người chơi nhận ra tiếng Anh chính là \textit{"Phép thuật"}. Việc nhập đúng từ vựng chính là gọi \textit{"Chân Danh"} (True Name) của sự vật để phá vỡ lớp bảo vệ của quái vật. Bối cảnh này là nền tảng để thiết kế giao diện và cơ chế chiến đấu của game.
\subsubsection*{b. Frontend (Giao diện \& Trải nghiệm Người dùng)}
Toàn bộ luồng giao diện được thiết kế sơ đồ hóa trên Figma trước khi lập trình trên Godot, chia làm hai trọng tâm chính:
\begin{itemize}
	\item \textbf{Quiz UI (\texttt{quiz\_ui.gd}):} Đảm nhiệm vai trò Placement Test (Kiểm tra đầu vào). Hệ thống cung cấp bộ câu hỏi tĩnh để đánh giá năng lực, sau đó gửi điểm số cho AI phân tích điểm yếu và xếp loại người chơi theo chuẩn CEFR. Giao diện này sẽ khóa mọi tương tác trong lúc chờ AI xử lý, đảm bảo an toàn dữ liệu.
    \item \textbf{Battle Scene (\texttt{battle\_scene.gd}):} Thiết kế theo dạng Combat Loop trực quan. Khi người chơi chạm trán quái vật, UI hiển thị các tùy chọn (Ma thuật) thông qua câu hỏi MCQ. Nếu người chơi trả lời sai, Pop-up "AI Tutor" (với nhân dạng Elaria) sẽ lập tức xuất hiện để giải thích lỗi sai, hiện thực hóa "Learning Loop" ngay trong trận đấu.
\end{itemize}

\subsubsection*{c. Backend (Hệ thống AI \& Thuật toán - \texttt{AIManager.gd})}
Backend của game đóng vai trò là bộ não xử lý logic học tập cá nhân hóa thông qua 3 kỹ thuật cốt lõi:
\begin{itemize}
    \item \textbf{Kiến trúc Nạp ngầm Bất đồng bộ (Asynchronous Pre-fetching):} Khắc phục nhược điểm tốc độ xử lý của Local LLM bằng cách thiết lập các hàng đợi đa tầng (Multi-tier Queues). Bằng các lệnh gọi API (\texttt{HTTPRequest}) liên tục qua cổng mạng 11434 của Ollama, hệ thống âm thầm chuẩn bị sẵn 5 câu hỏi cho mỗi khu vực (Biome) ngay khi người chơi đang di chuyển, đưa độ trễ tải câu hỏi trong trận chiến về mức 0 giây.
    \item \textbf{Dynamic Difficulty Adjustment (DDA - Thích nghi Độ khó):} Thuật toán DDA tính toán chỉ số thông thạo (\textit{mastery\_score}) của người chơi đối với từng từ vựng riêng biệt. Dựa trên chỉ số này, Backend tự động phân loại và chỉ thị AI sinh ra các dạng câu hỏi khác nhau (Nhận diện nghĩa cơ bản cho từ mới gặp; Điền vào chỗ trống cho từ hay sai; Phân biệt sắc thái đồng/trái nghĩa cho từ đã thành thạo).
    \item \textbf{Prompt Engineering:} Kích hoạt kỷ luật thép cho LLM thông qua các Prompt ép khuôn. Bắt buộc mô hình Qwen2.5:3b trả về chuỗi định dạng JSON chuẩn (chứa câu hỏi, 4 đáp án, logic đáp án đúng và giải thích tiếng Việt), đồng thời giới hạn \textit{temperature} để triệt tiêu hiện tượng AI ảo giác (Hallucination) hoặc sinh câu hỏi lạc đề.
\end{itemize}

\subsubsection*{d. Database (Cơ sở Dữ liệu - \texttt{DatabaseManager.gd})}
Hệ thống dữ liệu (Single Source of Truth) của dự án quản lý toàn bộ luồng thông tin người chơi:
\begin{itemize}
    \item \textbf{Kiến trúc Relational DB:} Dữ liệu được tổ chức thành 4 bảng quan hệ logic: \textit{Player\_Profile} (lưu vị trí, HP, CEFR level), \textit{Enemy\_Tier\_Dict} (cấp độ quái), \textit{Vocabulary\_Bank} (kho từ vựng), và \textit{Player\_Vocab\_Mastery} (lưu lịch sử tương tác đúng/sai của từng từ).
    \item \textbf{Data Pipeline:} Hệ thống có khả năng tự động Seed (nạp) hàng loạt từ vựng mới từ file \texttt{vocabulary.csv} bên ngoài thông qua một bộ Parser nội bộ, giúp dễ dàng mở rộng nội dung game mà không cần can thiệp vào mã nguồn.
    \item \textbf{Decisions \& Trade-offs (Quyết định thiết kế):} Nhóm ưu tiên sử dụng SQLite (Local Database) thay vì các giải pháp Cloud DB (như Firebase). Nguyên nhân cốt lõi là do Backend AI (Ollama) đang chạy hoàn toàn trên máy trạm của người dùng. Việc kết hợp với một Local DB sẽ giúp luồng dữ liệu (Data Flow) đồng nhất, loại bỏ hoàn toàn rủi ro rớt mạng (Network Latency), đảm bảo trò chơi có thể vận hành mượt mà $100\%$ Offline.
\end{itemize}